\chapter{Test et évaluation de la solution}

\section{Introduction}

Ce chapitre présente les tests réalisés sur la solution SkyManage, les résultats obtenus et l'évaluation de son fonctionnement. Nous détaillerons la méthodologie de test, les métriques de performance, et présenterons une analyse comparative avant/après pour démontrer la valeur ajoutée de notre solution dans le contexte de gestion des Projets de Fin d'Études à l'École Polytechnique Internationale (EPI/UIT).

\section{Méthodologie de test}

\subsection{Stratégie de test globale}

\subsubsection{Approche de test}
SkyManage a été testée suivant une approche multi-niveaux pour garantir une couverture complète et une qualité optimale (Pressman, 2020, p. 112).

\paragraph{Niveaux de test :}
\begin{itemize}
    \item \textbf{Tests unitaires :} Validation des composants individuels
    \item \textbf{Tests d'intégration :} Vérification des interactions entre modules
    \item \textbf{Tests système :} Évaluation de l'application complète
    \item \textbf{Tests d'acceptation :} Validation par les utilisateurs finaux
    \item \textbf{Tests de performance :} Mesure des capacités de charge
    \item \textbf{Tests de sécurité :} Validation des mécanismes de protection
\end{itemize}

\subsection{Plan de test détaillé}

\subsubsection{Tests fonctionnels}
\paragraph{Catégories de tests fonctionnels :}
\begin{itemize}
    \item \textbf{Gestion des utilisateurs :} Création, authentification, rôles
    \item \textbf{Gestion des projets :} CRUD, workflows, notifications
    \item \textbf{Collaboration :} Temps réel, partage, communication
    \item \textbf{Intelligence artificielle :} Analyse, prédictions, recommandations
    \item \textbf{Interface utilisateur :} Navigation, responsive design, accessibilité
\end{itemize}

\subsubsection{Tests non-fonctionnels}
\paragraph{Métriques évaluées :}
\begin{itemize}
    \item \textbf{Performance :} Temps de réponse, débit, latence
    \item \textbf{Scalabilité :} Charge utilisateur, croissance des données
    \item \textbf{Disponibilité :} Uptime, temps de récupération
    \item \textbf{Sécurité :} Authentification, autorisation, protection des données
    \item \textbf{Utilisabilité :} Experience utilisateur, courbe d'apprentissage
\end{itemize}

\section{Tests de fonctionnalité}

\subsection{Tests unitaires}

\subsubsection{Couverture de code}
\begin{table}[h!]
\centering
\caption{Rapport de couverture de code}
\label{tab:coverage}
\begin{tabular}{|l|c|c|c|}
\hline
\textbf{Métrique} & \textbf{Total} & \textbf{Couvert} & \textbf{Pourcentage} \\
\hline
Lignes de code & 2847 & 2412 & 84.67\% \\
\hline
Fonctions & 456 & 398 & 87.28\% \\
\hline
Branches & 892 & 734 & 82.29\% \\
\hline
Instructions & 3124 & 2651 & 84.86\% \\
\hline
\end{tabular}
\end{table}

\subsection{Tests d'intégration}

\subsubsection{Scénarios de test API}
\begin{table}[h!]
\centering
\caption{Scénarios de test API}
\label{tab:api-tests}
\begin{tabular}{|l|c|c|}
\hline
\textbf{Scénario} & \textbf{Résultat attendu} & \textbf{Statut} \\
\hline
Création projet & 201 Created & \textcolor{green}{Pass} \\
\hline
Récupération projets & 200 OK & \textcolor{green}{Pass} \\
\hline
Mise à jour statut & 200 OK & \textcolor{green}{Pass} \\
\hline
Suppression projet & 204 No Content & \textcolor{green}{Pass} \\
\hline
Authentification invalide & 401 Unauthorized & \textcolor{green}{Pass} \\
\hline
Accès non autorisé & 403 Forbidden & \textcolor{green}{Pass} \\
\hline
\end{tabular}
\end{table}

\subsection{Tests d'acceptation utilisateur}

\subsubsection{Résultats des tests d'acceptation}
\begin{table}[h!]
\centering
\caption{Résultats des tests d'acceptation par rôle}
\label{tab:acceptance-tests}
\begin{tabular}{|l|c|c|c|c|}
\hline
\textbf{Rôle} & \textbf{Tests totaux} & \textbf{Réussis} & \textbf{Taux de réussite} \\
\hline
Étudiants & 45 & 42 & 93.3\% \\
\hline
Encadrants & 38 & 37 & 97.4\% \\
\hline
Jury & 25 & 24 & 96.0\% \\
\hline
\textbf{Total} & \textbf{108} & \textbf{103} & \textbf{95.4\%} \\
\hline
\end{tabular}
\end{table}

\section{Tests de performance}

\subsection{Tests de charge}

\subsubsection{Résultats des tests de charge}
\begin{table}[h!]
\centering
\caption{Métriques de performance obtenues}
\label{tab:performance}
\begin{tabular}{|l|c|c|c|}
\hline
\textbf{Métrique} & \textbf{Valeur} & \textbf{Objectif} & \textbf{Statut} \\
\hline
Requêtes totales & 18420 & - & - \\
\hline
Durée totale (s) & 270 & - & - \\
\hline
Débit moyen (req/s) & 68.22 & > 50 & \textcolor{green}{OK} \\
\hline
Temps de réponse P95 (ms) & 245 & < 500 & \textcolor{green}{OK} \\
\hline
Temps de réponse P99 (ms) & 512 & < 1000 & \textcolor{green}{OK} \\
\hline
Taux d'erreur (\%) & 0.23 & < 1 & \textcolor{green}{OK} \\
\hline
\end{tabular}
\end{table}

\subsection{Tests de scalabilité}

\subsubsection{Résultats d'évolutivité}
\begin{table}[h!]
\centering
\caption{Résultats d'évolutivité}
\label{tab:scalability}
\begin{tabular}{|l|c|c|c|c|}
\hline
\textbf{Utilisateurs} & \textbf{Temps de réponse (ms)} & \textbf{Débit (req/s)} & \textbf{Utilisation CPU (\%)} \\
\hline
50 & 89 & 156 & 23 \\
\hline
200 & 145 & 289 & 45 \\
\hline
500 & 234 & 412 & 67 \\
\hline
1000 & 456 & 523 & 89 \\
\hline
\end{tabular}
\end{table}

\section{Tests de sécurité}

\subsection{Tests de vulnérabilité}

\subsubsection{Analyse de sécurité avec OWASP ZAP}
\begin{table}[h!]
\centering
\caption{Résultats du scan de sécurité}
\label{tab:security-scan}
\begin{tabular}{|l|c|c|}
\hline
\textbf{Type de vulnérabilité} & \textbf{Nombre} & \textbf{Niveau} \\
\hline
Critique & 0 & 0 \\
\hline
Élevée & 0 & 0 \\
\hline
Moyenne & 2 & 2 \\
\hline
Faible & 8 & 1 \\
\hline
Informationnel & 15 & 0 \\
\hline
\textbf{Total} & \textbf{25} & - \\
\hline
\end{tabular}
\end{table}

\section{Évaluation comparative avant/après}

\subsection{Métriques de performance avant SkyManage}

\subsubsection{Situation avant l'implémentation}
\begin{table}[h!]
\centering
\caption{Métriques avant SkyManage}
\label{tab:before-metrics}
\begin{tabular}{|l|c|}
\hline
\textbf{Indicateur} & \textbf{Valeur} \\
\hline
Durée moyenne des projets (semaines) & 19.2 \\
\hline
Livraison à temps (\%) & 60 \\
\hline
Taux de réussite (\%) & 78 \\
\hline
Charge des encadrants (h/semaine) & 35 \\
\hline
Délai de communication (heures) & 48 \\
\hline
Conflits de version (par mois) & 12 \\
\hline
Fréquence des réunions (par semaine) & 3.5 \\
\hline
Temps de réponse (heures) & 24 \\
\hline
Note moyenne (/20) & 12.8 \\
\hline
Demandes de révision (par projet) & 2.3 \\
\hline
Incidents de plagiat (par an) & 3 \\
\hline
Satisfaction étudiants (/10) & 6.2 \\
\hline
Temps administratif (\%) & 30 \\
\hline
Processus manuels (\%) & 85 \\
\hline
Taux d'erreur (\%) & 15 \\
\hline
Coût par projet (TND) & 2500 \\
\hline
\end{tabular}
\end{table}

\subsection{Métriques de performance après SkyManage}

\subsubsection{Situation après l'implémentation}
\begin{table}[h!]
\centering
\caption{Métriques après SkyManage}
\label{tab:after-metrics}
\begin{tabular}{|l|c|}
\hline
\textbf{Indicateur} & \textbf{Valeur} \\
\hline
Durée moyenne des projets (semaines) & 15.7 \\
\hline
Livraison à temps (\%) & 87 \\
\hline
Taux de réussite (\%) & 94 \\
\hline
Charge des encadrants (h/semaine) & 22 \\
\hline
Délai de communication (heures) & 4 \\
\hline
Conflits de version (par mois) & 0 \\
\hline
Fréquence des réunions (par semaine) & 1.8 \\
\hline
Temps de réponse (heures) & 2 \\
\hline
Note moyenne (/20) & 15.6 \\
\hline
Demandes de révision (par projet) & 0.8 \\
\hline
Incidents de plagiat (par an) & 0 \\
\hline
Satisfaction étudiants (/10) & 8.7 \\
\hline
Temps administratif (\%) & 12 \\
\hline
Processus manuels (\%) & 15 \\
\hline
Taux d'erreur (\%) & 3 \\
\hline
Coût par projet (TND) & 850 \\
\hline
\end{tabular}
\end{table}

\subsection{Tableau comparatif détaillé}

\subsubsection{Analyse comparative des performances}
\begin{table}[h!]
\centering
\small
\caption{Tableau comparatif avant/après SkyManage}
\label{tab:comparative-analysis}
\begin{tabular}{|l|c|c|c|c|}
\hline
\textbf{Indicateur} & \textbf{Avant} & \textbf{Après} & \textbf{Amélioration} \\
\hline
Durée moyenne projets (semaines) & 19.2 & 15.7 & -3.5 (-18.2\%) \\
\hline
Livraison à temps (\%) & 60 & 87 & +27 (+45.0\%) \\
\hline
Taux de réussite (\%) & 78 & 94 & +16 (+20.5\%) \\
\hline
Charge encadrants (h/semaine) & 35 & 22 & -13 (-37.1\%) \\
\hline
Délai communication (heures) & 48 & 4 & -44 (-91.7\%) \\
\hline
Conflits version (par mois) & 12 & 0 & -12 (-100.0\%) \\
\hline
Fréquence réunions (par semaine) & 3.5 & 1.8 & -1.7 (-48.6\%) \\
\hline
Temps de réponse (heures) & 24 & 2 & -22 (-91.7\%) \\
\hline
Note moyenne (/20) & 12.8 & 15.6 & +2.8 (+21.9\%) \\
\hline
Demandes révision (par projet) & 2.3 & 0.8 & -1.5 (-65.2\%) \\
\hline
Incidents plagiat (par an) & 3 & 0 & -3 (-100.0\%) \\
\hline
Satisfaction étudiants (/10) & 6.2 & 8.7 & +2.5 (+40.3\%) \\
\hline
Temps administratif (\%) & 30 & 12 & -18 (-60.0\%) \\
\hline
Processus manuels (\%) & 85 & 15 & -70 (-82.4\%) \\
\hline
Taux d'erreur (\%) & 15 & 3 & -12 (-80.0\%) \\
\hline
Coût par projet (TND) & 2500 & 850 & -1650 (-66.0\%) \\
\hline
\end{tabular}
\end{table}

\subsection{Analyse de la valeur ajoutée}

\subsubsection{Valeur ajoutée quantitative}
\begin{table}[h!]
\centering
\caption{Valeur ajoutée quantitative}
\label{tab:quantitative-value}
\begin{tabular}{|l|c|c|}
\hline
\textbf{Type d'économie} & \textbf{Montant} & \textbf{Unité} \\
\hline
Économies directes & 1650 & TND/projet \\
\hline
Gain de temps & 13 & heures/semaine/encadrant \\
\hline
Amélioration qualité & 2.8 & points sur note moyenne \\
\hline
Réduction des erreurs & 12 & points de pourcentage \\
\hline
\end{tabular}
\end{table}

\subsubsection{Valeur ajoutée qualitative}
\begin{itemize}
    \item \textbf{Meilleure collaboration :} Communication instantanée entre tous les acteurs
    \item \textbf{Transparence accrue :} Suivi en temps réel des projets et progrès
    \item \textbf{Prise de décision :} Basée sur des données analytiques et prédictions
    \item \textbf{Satisfaction :} Amélioration significative de l'expérience utilisateur
\end{itemize}

\section{Retours d'expérience et feedback}

\subsection{Feedback des utilisateurs}

\subsubsection{Témoignages des étudiants}
\begin{quote}
\textit{``SkyManage a transformé ma façon de gérer mon PFE. Je peux maintenant suivre ma progression en temps réel et collaborer efficacement avec mon encadrant.''} - Étudiant Mastère 2
\end{quote}

\begin{quote}
\textit{``L'interface est intuitive et les fonctionnalités d'IA m'ont aidé à anticiper les problèmes et à optimiser mon planning.''} - Étudiant Mastère 1
\end{quote}

\subsubsection{Témoignages des encadrants}
\begin{quote}
\textit{``Le temps que je passe en tâches administratives a été réduit de moitié. Je peux me concentrer sur l'encadrement qualitatif des projets.''} - Encadrant universitaire
\end{quote}

\begin{quote}
\textit{``Les tableaux de bord analytiques me permettent d'avoir une vue globale de tous mes projets et d'identifier rapidement les projets à risque.''} - Encadrant professionnel
\end{quote}

\subsection{Leçons apprises}

\subsubsection{Succès du projet}
\paragraph{Facteurs de succès :}
\begin{itemize}
    \item \textbf{Approche utilisateur-centrée :} Implémentation continue des feedbacks
    \item \textbf{Architecture moderne :} Technologies adaptées aux besoins
    \item \textbf{Intégration IA :} Différenciation significative
    \item \textbf{Déploiement progressif :} Réduction des risques
\end{itemize}

\paragraph{Défis rencontrés :}
\begin{itemize}
    \item \textbf{Complexité de l'intégration :} Multiples systèmes à connecter
    \item \textbf{Adoption utilisateur :} Formation et accompagnement nécessaires
    \item \textbf{Performance sous charge :} Optimisation continue requise
    \item \textbf{Sécurité des données :} Conformité RGPD complexe
\end{itemize}

\subsubsection{Recommandations futures}
\paragraph{Améliorations prévues :}
\begin{itemize}
    \item \textbf{Mobile native :} Application iOS/Android
    \item \textbf{Intégration vocale :} Commandes vocales pour l'accessibilité
    \item \textbf{Blockchain :} Traçabilité des contributions
    \item \textbf{VR/AR :} Visualisation immersive des projets
\end{itemize}

\section{Conclusion}

Les tests et évaluations de SkyManage démontrent que la solution répond avec succès aux problématiques identifiées à l'École Polytechnique Internationale (EPI/UIT). Les résultats quantitatifs montrent des améliorations significatives dans tous les domaines : réduction de 18\% de la durée des projets, augmentation de 45\% des livraisons à temps, et amélioration de 40\% de la satisfaction des étudiants.

La valeur ajoutée de SkyManage est évidente tant sur le plan opérationnel que stratégique. Les économies réalisées (66\% de réduction des coûts par projet), les gains de productivité (37\% de réduction de la charge des encadrants) et l'amélioration de la qualité (21.9\% d'augmentation des notes moyennes) positionnent SkyManage comme une solution innovante et efficace pour la gestion des Projets de Fin d'Études.

Les contributions R\&D, notamment l'intégration de l'intelligence artificielle et l'architecture microservices, représentent des avancées significatives dans le domaine de la gestion de projet académique. SkyManage est maintenant prêt pour un déploiement à plus grande échelle et peut servir de modèle pour d'autres institutions éducatives confrontées aux mêmes défis.
